%!TEX root = ../template.tex
%%%%%%%%%%%%%%%%%%%%%%%%%%%%%%%%%%%%%%%%%%%%%%%%%%%%%%%%%%%%%%%%%%%%
%% glossary.tex
%% NOVA thesis document file
%%
%% Glossary definition
%%%%%%%%%%%%%%%%%%%%%%%%%%%%%%%%%%%%%%%%%%%%%%%%%%%%%%%%%%%%%%%%%%%%

\typeout{NT FILE glossary.tex}%

\newglossaryentry{computer}{
  name={computer},
  sort={computer},
  description={An electronic device which is capable of receiving information (data) in a particular form and of performing a sequence of operations in accordance with a predetermined but variable set of procedural instructions (program) to produce a result in the form of information or signals. This is a test that adds a citation~\cite{Artho04} to the glossary!}
}
\newglossaryentry{uml}{
  name={UML},
  description={Unified Modeling Language, a standardized modeling language in software engineering used to specify, visualize, develop, and document software system artifacts},
}

\newglossaryentry{generativeai}{
  name={Generative AI},
  description={Artificial intelligence models that generate content such as text, images, and code, often using large-scale neural networks like GPT or DALL-E},
}

\newglossaryentry{requirementmodeling}{
  name={Requirement Modeling},
  description={The process of defining, analyzing, and documenting software requirements to create structured representations of system needs},
}

\newacronym{llm}{LLM}{Large Language Model}
\newacronym{ner}{NER}{Named Entity Recognition}
\newacronym{cot}{CoT}{Chain of Thought}
\newacronym{rouge}{ROUGE}{Recall-Oriented Understudy for Gisting Evaluation}
\newacronym {srs}{SRS}{Software Requirements Specification}
